\section{Main Experiment Results}

Table~\ref{tab:main-results} summarizes the current three-model comparison. Exact values for best accuracy, completed epochs, and average epoch times should be copied from the finalized Weights \& Biases run summaries before submission. They are intentionally left as explicit placeholders in this draft to avoid estimating quantitative values from screenshots.

\begin{table*}[t]
    \centering
    \small
    \caption{Main CIFAR-10 comparison for the current reproduction pipeline. Parameter counts are exact for the present model definitions; the remaining metrics should be copied from the corresponding finalized Weights \& Biases runs.}
    \label{tab:main-results}
    \begin{tabular}{lcccccc}
        \toprule
        Model & Params & Epochs & Best acc. & Train s/epoch & Eval s/epoch & Notes \\
        \midrule
        ConvMixer & \convmixerparams & \convmixerepochs & \convmixeracc & \convmixertrainsec & \convmixerevalsec & best acc., slowest \\
        ResNet18 & \resnetparams & \resnetepochs & \resnetacc & \resnettrainsec & \resnetevalsec & strong CNN baseline \\
        DeiT-Tiny & \deitparams & \deitepochs & \deitacc & \deittrainsec & \deitevalsec & weakest current run \\
        \bottomrule
    \end{tabular}
\end{table*}

\begin{figure*}[t]
    \centering
    \IfFileExists{figures/main_accuracy.pdf}{
        \includegraphics[width=0.95\textwidth]{figures/main_accuracy.pdf}
    }{
        \IfFileExists{figures/main_accuracy.png}{
            \includegraphics[width=0.95\textwidth]{figures/main_accuracy.png}
        }{
            \fbox{
                \parbox[c][2.1in][c]{0.92\textwidth}{
                    \centering
                    Export the Weights \& Biases plot for \texttt{best\_test\_accuracy} vs. epoch and save it as\\
                    \texttt{report/figures/main\_accuracy.pdf} or \texttt{report/figures/main\_accuracy.png}.\\[0.6em]
                    The fallback screenshot should include only the three final runs: ConvMixer, ResNet18, and DeiT-Tiny.
                }
            }
        }
    }
    \caption{Main training comparison from Weights \& Biases. The preferred export is the \texttt{best\_test\_accuracy} plot for the three final runs.}
    \label{fig:main-accuracy}
\end{figure*}

The qualitative ordering is already clear from the current logged runs. ConvMixer achieves the highest test accuracy in the present setup and reaches its plateau relatively quickly. ResNet18 remains competitive and substantially stronger than DeiT-Tiny, but it does not reach the same top accuracy as ConvMixer. DeiT-Tiny improves steadily, yet under the current training configuration it lags well behind the other two models on CIFAR-10.

Runtime behavior tells a different story. The same Weights \& Biases workspace shows that ConvMixer is slower per training epoch than both ResNet18 and DeiT-Tiny in the current runs, even though it is much smaller than ResNet18 in parameter count. This suggests that the depthwise-plus-pointwise structure is not automatically the cheapest option in wall-clock terms under the present implementation and execution environment.

The result pattern is also consistent with the intended interpretation of the benchmark. ResNet18 benefits from a strong convolutional prior and remains a reliable baseline on small natural-image datasets. DeiT-Tiny, in contrast, is more sensitive to recipe details and may need a more specialized optimization setup to close the gap. ConvMixer occupies an interesting middle ground: it keeps explicit convolutional mixing while adopting patch-based processing, and in this reproduction that combination appears to be the most effective among the tested options.

This comparison should still be interpreted with care. The current runs are useful for establishing a ranking inside one codebase, but they are not necessarily a perfectly matched benchmark. If the final ConvMixer, ResNet18, and DeiT-Tiny runs differ in epoch budget, scheduler horizon, or runtime backend, then the comparison remains valid as an empirical observation inside the reproduction pipeline but not as a fully controlled claim about architecture quality in general.

\begin{figure*}[t]
    \centering
    \IfFileExists{figures/main_runtime.pdf}{
        \includegraphics[width=0.95\textwidth]{figures/main_runtime.pdf}
    }{
        \IfFileExists{figures/main_runtime.png}{
            \includegraphics[width=0.95\textwidth]{figures/main_runtime.png}
        }{
            \fbox{
                \parbox[c][1.7in][c]{0.92\textwidth}{
                    \centering
                    Optional runtime figure. If space permits, export the Weights \& Biases comparison for\\
                    \texttt{time/train\_sec} and \texttt{time/eval\_sec} and save it as\\
                    \texttt{report/figures/main\_runtime.pdf} or \texttt{report/figures/main\_runtime.png}.
                }
            }
        }
    }
    \caption{Optional runtime comparison from Weights \& Biases. This figure is useful if the final draft still has room after the main accuracy plot and result table.}
    \label{fig:main-runtime}
\end{figure*}
